\section{Open-World Mechanism Discovery}

A protected canonical state can still be incomplete in a systematic way. The failure is not necessarily shallow search. A recursive process can generate thousands of descendants while remaining inside the vocabulary, citation neighborhood and disciplinary assumptions of its first query. This is \emph{ontology-conditioned closure}: depth increases while the accessible concept space does not.

Information retrieval has long documented the vocabulary problem---different people use different words for the same objects and functions \cite{furnas1987}. Literature-based discovery showed that scientifically useful connections can exist between literatures that do not directly cite one another \cite{swanson1986}. More recent methodology-inspiration retrieval explicitly targets transferable methods rather than ordinary topical similarity \cite{garikaparthi2025}. Search diversification likewise treats breadth across intents or aspects as an objective distinct from relevance alone \cite{carbonell1998}. OWMD turns these observations into a governance requirement for research-agent architecture.

\subsection{Capability ownership before search}

For a subsystem \(M\), let
\[
\mathcal F_{\mathrm{req}}(M)=\{f_1,\ldots,f_n\}
\]
be the functions that the subsystem must perform. A function is considered \emph{owned} only when there exists a record
\[
O_f=(m_f,\sigma_f,\mathrm{pre}_f,\mathrm{post}_f,E_f,T_f,\mathrm{fail}_f),
\]
containing a mechanism identity \(m_f\), scope \(\sigma_f\), preconditions, postconditions, evidence, executable tests and failure semantics. A label such as ``memory,'' ``reasoning'' or ``workspace'' is not an owner record. This requirement follows the spirit of explicit contracts in software and formal methods: the question is what the mechanism guarantees under stated conditions, not what it is called.

For every unowned or contested function, construct a functional signature
\[
\Sigma(f)=(I,O,C,R,D,X),
\]
where \(I\) and \(O\) are characteristic inputs and outputs, \(C\) resource or structural constraints, \(R\) relations among participating objects, \(D\) temporal/control dynamics, and \(X\) failure or intervention signatures. The signature is written so that current framework labels can be masked.

For the bounded-selection/broadcast function, for example, a signature can say: many processes generate candidates; a small active subset is selected competitively; capacity is limited; selected objects become reusable by heterogeneous downstream processes; items can persist, be displaced or be evicted; interventions on active contents can change later computation; and computational prominence does not establish truth. None of those phrases requires the terms ``global workspace,'' ``blackboard'' or ``J-space.''

\subsection{Route ensemble}

OWMD executes a heterogeneous route family \(\mathcal R\). The reference family includes:

\begin{enumerate}[leftmargin=*]
\item exact terminology;
\item lexical variants and paraphrases;
\item a function-only route with current core vocabulary masked;
\item historical precursors;
\item mathematical equivalents or neighboring formalisms;
\item implementation analogues from other technical communities;
\item methodology-inspiration retrieval;
\item backward and forward citation-neighborhood expansion;
\item literature bridges between otherwise separated topics;
\item adversarial alternatives that solve the function differently;
\item cross-language expansion when the domain makes it relevant; and
\item a final freshness scan reaching a declared cutoff.
\end{enumerate}

Every route retains its query, route kind, completion state, candidate identities, lexical-independence flag, stability status and freshness metadata. At least one completed route must be lexically independent of the current core vocabulary. Lexical independence is tested on the actual query tokens rather than accepted as a declaration.

The goal is not to maximize query count. Routes are designed to fail differently. Exact terminology has high precision after a name is known but cannot recover an unknown name. Function-only search is robust to nomenclature but can be broad. Historical search can find mechanisms whose modern vocabulary changed. Mathematical-equivalence routes can reveal an older formal result hiding behind different application language. Implementation analogues can expose engineering solutions that are absent from the target discipline. Citation expansion exploits local graph structure but risks remaining inside one community. The ensemble therefore treats route diversity as an epistemic defense.

Open-world learning and open-set recognition address a different but related problem: systems must detect and manage classes that were not fully represented at training time \cite{jafarzadeh2021}. OWMD is not a classifier for unknown classes. Its common principle is the refusal to assume that the current ontology exhausts the world.

\subsection{Candidate assimilation and novelty discipline}

A retrieved mechanism \(m\) is not simply appended as a citation. It is assimilated relative to the target function with a status
\[
\mathrm{assim}(m,f)\in
\left\{
\begin{array}{lll}
\mathrm{EQUIVALENT}, & \mathrm{SUBSUMED}, & \mathrm{COMPLEMENTARY},\\
\mathrm{CONFLICTING}, & \mathrm{NOVEL\ RESIDUAL}, & \mathrm{UNRESOLVED}
\end{array}
\right\}.
\]
The candidate record retains source identifiers, route identifiers, functional and structural fit, evidence quality, novelty threat, transfer cost and contradiction flags.

This step protects the novelty claim. If a historical system already implements the function under materially equivalent assumptions, the correct action is to narrow the RAKL contribution. If the old system solves a strict superset of the function, the new mechanism may be subsumed. If it contributes a complementary property---for example, strong computational access but no evidence governance---the relationship should be stated as a decomposition rather than a priority claim. Unresolved candidates remain explicit research fibers.

\subsection{Discovery workspace}

Retrieval can still be locally myopic if the candidate shortlist is ranked only by current relevance. OWMD therefore has its own bounded Discovery Workspace with partitions
\[
\{\mathrm{NEAR},\mathrm{REMOTE},\mathrm{CHALLENGE},\mathrm{HISTORICAL},\mathrm{FRESH}\}.
\]
The reference gate reserves slots for remote, challenging, historical and fresh candidates before filling by global priority. A missing reserved partition fails closed. The semantics are analogous to the research workspace of Section~6, but the object being diversified is a mechanism candidate rather than an active scientific claim.

\subsection{Bounded discovery closure}

Let \(Q\) denote the discovery question, \(\mathcal D\) a declared source universe, \(B\) finite route/retrieval budgets and \(t^\star\) a freshness cutoff.

\begin{definition}[Bounded discovery closure]
A required function \(f\) is \emph{bounded-discovery-closed} relative to \((Q,\mathcal R,\mathcal D,B,t^\star)\) when:
\begin{enumerate}[leftmargin=*]
\item it has a complete owner or an explicit unresolved research fiber;
\item every mandatory route is complete or explicitly inapplicable;
\item at least one completed route is lexically independent of the current core vocabulary;
\item citation-neighborhood expansion satisfies its declared stability rule;
\item the freshness route reaches \(t^\star\);
\item an omission audit and nearest-work equivalence audit have provenance-bearing records; and
\item every unresolved retrieved candidate is preserved as an explicit fiber.
\end{enumerate}
\end{definition}

The reference implementation goes further than a boolean checklist: omission and nearest-work audits carry identities, reviewer-context provenance and evidence references. This avoids a degenerate closure call in which the same code that wants to stop simply passes \texttt{True} for ``independent review complete.''

\begin{theorem}[Unrestricted open-world completeness is not finitely certifiable]
Assume the concept/source universe is unrestricted and no complete membership oracle is available. No finite OWMD transcript can certify that no relevant undiscovered mechanism exists outside the retrieved set.
\end{theorem}

\begin{proof}
A finite run issues finitely many queries and observes finitely many returned objects. Consider two possible worlds that agree on the complete observed transcript. In the first world, no additional relevant mechanism exists. In the second, add one relevant mechanism that is described only in a vocabulary or source not returned by any executed route. Because the observed transcript is identical in both worlds, no decision function over that transcript can distinguish them. A universal non-existence claim is therefore not identifiable from the finite transcript.
\end{proof}

This result is the formal reason the software has no \texttt{ABSOLUTELY\_COMPLETE} success state. It can report only \texttt{BOUNDED\_CLOSED} relative to a declared search world.

\subsection{GWT-OMISSION-01}

The development history supplies a regression test. The benchmark withholds the names ``global workspace,'' ``consciousness,'' ``J-space,'' ``blackboard'' and relevant author names. It exposes only the functional signature described above. A successful scored retrieval must recover mechanism families containing blackboard architectures, Global Workspace/Global Neuronal Workspace, the Consciousness Prior, shared neural workspaces and the J-space study through at least one lexically independent route.

The benchmark currently validates the scoring, provenance and name-leakage contract. It does not establish prospective recall on arbitrary future hidden concepts. That stronger statement would require fresh pre-registered concepts whose identities are genuinely unavailable to the search process before execution.

\section{Bounded epistemic saturation}

OWMD says when one required function has been searched broadly enough to receive a bounded closure certificate. The project-level stopping rule is larger. A paper or research program can still grow after its mechanism owners are known: a new derivation may appear, an assumption may need narrowing, a counterexample may invalidate a theorem, or a newly independent experiment may change an evidence coordinate. The system therefore needs a definition of \emph{epistemic saturation} rather than search completion alone.

The term ``saturation'' has several established meanings. In qualitative research it commonly refers to conditions under which further sampling or analysis ceases to add relevant conceptual material, and the literature stresses that saturation must be operationalized relative to the research question and analytic framework \cite{saunders2018}. In automated reasoning and program optimization, saturation procedures repeatedly apply inference or rewrite rules until no new consequences or equivalences are produced; equality saturation uses e-graphs to preserve many equivalent expressions while rewrites accumulate \cite{tate2009,willsey2021}. Database and logic-programming fixed-point semantics likewise compute consequences by iterating monotone operators \cite{vanemden1976}. These traditions motivate the fixed-point language, but none by itself provides the evidence, novelty, freshness and open-world conditions required here.

\subsection{Saturation basis}

A flat growth trace is meaningful only if the measurement semantics stay fixed. Define a saturation basis
\[
\mathcal B=(b,\Omega,\iota,\rho,\nu,\epsilon),
\]
where \(b\) is a basis identity, \(\Omega\) the research/manuscript scope, \(\iota\) an object-identity policy, \(\rho\) a discovery-route-family version, \(\nu\) a novelty/equivalence policy and \(\epsilon\) an evidence/authority policy. The implementation fingerprints these coordinates content-wise. A change of basis invalidates a longitudinal saturation comparison rather than being counted as new knowledge.

This mirrors the metrology discipline used elsewhere in RAKL. Renaming or splitting a research fiber can change graph geometry without changing scientific content. Saturation therefore counts epistemic events, not raw representation edits.

\subsection{Marginal epistemic growth vector}

For research round \(r\), define
\[
\Delta_r=
(\Delta M,\Delta D,\Delta E,\Delta C,\Delta N,
\Delta V,\Delta A,\Delta F,\Delta R)_r,
\]
with non-negative coordinates:
\begin{align*}
\Delta M&=\text{new mechanisms or function owners},\\
\Delta D&=\text{new derivations or formal consequences},\\
\Delta E&=\text{new independent evidence roots},\\
\Delta C&=\text{new contradictions or counterexamples},\\
\Delta N&=\text{new negative results},\\
\Delta V&=\text{novelty/prior-art boundary updates},\\
\Delta A&=\text{assumption or scope updates},\\
\Delta F&=\text{unresolved-fiber updates},\\
\Delta R&=\text{discovery-route updates}.
\end{align*}
A round is substantively flat when \(\Delta_r=\mathbf 0\). Representation-only changes---sentence polishing, file reorganization, equivalent notation or ontology renaming under the same identity policy---are tracked separately and do not reset epistemic flatness.

The vector is deliberately non-compensatory. Ten new prose paragraphs cannot cancel one new counterexample; a large number of citations that duplicate existing evidence roots cannot compensate for an unresolved mechanism gap. This follows the same rationale as the earlier separation between description length, compression and target-relevant retained function. Minimum-description-length thinking warns against representational bookkeeping as progress \cite{rissanen1978}; the information bottleneck distinguishes compression from retained task-relevant information \cite{tishby2000}. RAKL uses these as analogies rather than claiming to optimize either formal objective.

\subsection{Finite-basis fixed point}

The fixed-point interpretation can be made exact in a bounded world. Let \(U_{\mathcal B}\) be a finite universe of distinguishable epistemic objects under basis \(\mathcal B\), and let
\[
F_{\mathcal B}:\powerset(U_{\mathcal B})\to\powerset(U_{\mathcal B})
\]
be a monotone, inflationary expansion operator representing all inference, retrieval and assimilation operations admitted by the basis. Starting from \(K_0\subseteq U_{\mathcal B}\), define
\[
K_{n+1}=F_{\mathcal B}(K_n).
\]

\begin{theorem}[Finite-basis saturation]
If \(U_{\mathcal B}\) is finite and \(F_{\mathcal B}\) is monotone and inflationary, then the iteration reaches a fixed point after at most \(|U_{\mathcal B}|-|K_0|\) strict-growth steps. The stabilized state is the least fixed point of \(F_{\mathcal B}\) containing \(K_0\).
\end{theorem}

\begin{proof}
Inflationarity gives \(K_n\subseteq K_{n+1}\). Every strict inclusion adds at least one element of the finite universe, so there can be at most \(|U_{\mathcal B}|-|K_0|\) strict-growth steps. At the first non-strict step, \(K_{n+1}=K_n\), hence \(K_n\) is a fixed point.

Now let \(Y\) be any fixed point with \(K_0\subseteq Y\). If \(K_n\subseteq Y\), monotonicity gives
\[
K_{n+1}=F_{\mathcal B}(K_n)\subseteq F_{\mathcal B}(Y)=Y.
\]
By induction every iterate is contained in \(Y\), including the stabilized state. Thus the stabilized state is the least fixed point above \(K_0\).
\end{proof}

This is a standard closure/fixed-point pattern closely related to Tarski's theorem \cite{tarski1955}. Its scientific value lies in exposing the assumptions that fail in an unrestricted open world. The real literature is not a fixed finite \(U_{\mathcal B}\); the source universe changes with time, languages, databases and newly discovered vocabularies. A paper can therefore be a fixed point only relative to its bounded expansion operator.

\subsection{Operational saturation certificate}

The software uses a stricter operational rule than ``one call to \(F\) added nothing.'' A bounded saturation report requires \(m\ge1\) consecutive flat rounds under the same basis fingerprint. For each of those final rounds:

\begin{enumerate}[leftmargin=*]
\item all required functions are bounded-discovery-closed;
\item route coverage is stable;
\item the omission audit passes;
\item the nearest-work equivalence audit passes;
\item no blocking research fibers remain; and
\item the freshness scan reaches the required cutoff.
\end{enumerate}

The reference default used for manuscript release is multiple consecutive flat rounds rather than a single flat round. This is not a statistical proof of recall. It is an adversarial engineering safeguard: the system must fail to grow after more than one independently designed expansion/review pass.

\begin{proposition}[New substantive knowledge reopens saturation]
If a bounded-saturated state is followed by a round \(r\) with \(\Delta_r\ne\mathbf0\), then the new state is not saturated under the consecutive-flat-round rule until the required number of subsequent flat rounds has again been observed.
\end{proposition}

\begin{proof}
The newest suffix of consecutive flat rounds has length zero when \(\Delta_r\ne\mathbf0\). Therefore it cannot satisfy any positive required flat-round count.
\end{proof}

This proposition captures the project principle directly: saturation is not a permanent badge. A new theorem, counterexample, primary source, prior-art equivalence or mechanism owner is evidence that the previous state was incomplete relative to what is now known, so the release gate reopens.

\subsection{Freshness expiry}

A saturation certificate also ages. Let \(t^\star\) be the freshness cutoff recorded by the final route scan and suppose a new release policy requires coverage through \(t'>t^\star\).

\begin{proposition}[Freshness-horizon expiry]
A bounded saturation certificate whose freshness cutoff is \(t^\star\) does not satisfy a release basis requiring cutoff \(t'>t^\star\) until the freshness route is rerun through at least \(t'\).
\end{proposition}

\begin{proof}
Freshness coverage through \(t'\) is an explicit predicate of the new basis. The old certificate establishes that predicate only through \(t^\star\), so the new obligation is unproved.
\end{proof}

This is why a paper can be saturated today and reopened by tomorrow's literature. Recent work on decision-theoretic stopping rules for document screening likewise treats stopping as dependent on the objective and the costs/benefits of further search rather than as an intrinsic property of a ranked list \cite{fletcher2026}. Systematic-review methodology similarly emphasizes explicit search strategies and update procedures rather than intuition about having seen enough papers \cite{macfarlane2022}.

\subsection{The paper as a bounded-saturated epistemic projection}

The long-form manuscript is generated under exactly this rule. Its state includes formal claims, proofs, counterexamples, evidence and literature links. A literature pass that discovers an equivalent mechanism changes \(\Delta V\); a mathematical audit that finds an omitted obstruction changes \(\Delta C\) or \(\Delta A\); an additional derivation changes \(\Delta D\); a new primary source can change \(\Delta E\) or \(\Delta V\). The manuscript must then grow or narrow accordingly.

Only when heterogeneous discovery/review passes become substantively flat under a fixed basis may the paper be described as \emph{bounded-saturated}. In the closure-system interpretation of Section~4, the manuscript describes one closed epistemic state---a fixed point relative to the declared operator---inside a genuine lattice of closed states. That is the precise sense in which the paper can be a description of a saturated epistemic-mechanics lattice. It is not a statement that the unrestricted scientific world has stopped growing.
